\cvsection{About me}

\begin{flushleft}
  \vspace{-2.0mm}
%  \entrydatestyle{
    % I was born on the 14th of January 1995 in Matera, Italy, where I lived and studied until 2013, the year in which I graduated at the Dante Alighieri High School of Science. From the beginning I was very interested in science and philosophy.
    % I moved to Milan, where I currently live, to study at the Politecnico di Milano in September of the same year. There I obtained a bachelor's degree and a master's degree in computer engineering, acquiring knowledge in the fields of calculus, statistics and computer science. I specialised in the discipline of artificial intelligence and became passionate about it during these years.
    % With a growing curiosity, I quickly immersed myself in its most applied aspects, focusing first on traditional machine learning and later on deep learning.
    % I am currently pursuing a PhD in Human-Centred artificial intelligence at the Politecnico di Milano. My research focuses on generative deep learning models and auto machine learning, with a particular emphasis on neural architecture search techniques. At the same time, I am interested in the application of these techniques in real-world contexts related to the clinical world, as well as in explainable artificial intelligence techniques and privacy in machine learning, with a particular focus on adversarial learning techniques.
    % My involvement includes research and teaching activities in academia and industry, and in both cases I use a scientific and ethical approach to guide my actions.
Eugenio Lomurno is a Postdoctoral Fellow at the Artificial Intelligence and Robotics Laboratory (AIRLab) of Politecnico di Milano, where he completed his entire academic journey, earning his Bachelor's degree (2017), Master's degree (2020), and PhD Cum Laude (2025) in Information Technology Engineering.\\
His educational background, rooted in solid scientific-philosophical foundations, has evolved towards complete specialisation in Artificial Intelligence, both theoretical and applied, a field that now represents the core of his research and teaching activities.
His current research focuses on developing synthetic datasets with high utility and privacy through advanced multimodal generative deep learning models, intended for training future neural architectures. His interests span from computer security applied to generative models - with particular focus on privacy in federated contexts - to neural architecture search, and the application of state-of-the-art techniques in oncology.\\
Within the scientific debate, he shares the concerns of that part of the scientific community which questions the possible future scenarios of artificial intelligence, advocating the necessity to develop strategies to address critical phenomena such as goal misalignment, self-replication, and uncontrolled self-improvement in advanced AI systems. 
\end{flushleft}
